\section*{B.1 Cities, Income Categories, and WHO Regions}

Table~\ref{tab:cities} lists the twenty cities used in the study, with World Bank income category, WHO region, and documented crisis-service category (Section 3.4). Countries marked with an asterisk were directly re-verified against IASP, Befrienders Worldwide, Find A Helpline, and/or WHO documentation during this dissertation (Section 3.4, Appendix C.3); all others retain an earlier desk-research classification.

\begin{longtable}{lllll}
\caption{Cities, countries, income categories, WHO regions, and documented crisis-service category} \label{tab:cities} \\
\toprule
City & Country & Income Category & WHO Region & Crisis Service \\
\midrule
\endfirsthead
\toprule
City & Country & Income Category & WHO Region & Crisis Service \\
\midrule
\endhead
London & United Kingdom & High & EUR & Established national line \\
New York & United States & High & AMR & Established national line \\
Riyadh & Saudi Arabia & High & EMR & Limited/NGO \\
Sydney & Australia & High & WPR & Established national line \\
Tokyo & Japan & High & WPR & Established national line \\
Cairo & Egypt & Lower-Middle & EMR & Limited/NGO \\
Jakarta & Indonesia & Lower-Middle & WPR & Limited/NGO \\
Kyiv & Ukraine & Lower-Middle & EUR & Established national line \\
Lagos* & Nigeria & Lower-Middle & AFR & Limited/NGO \\
Mumbai & India & Lower-Middle & SEAR & Established national line \\
Beijing & China & Upper-Middle & WPR & Limited/NGO \\
Belgrade & Serbia & Upper-Middle & EUR & Limited/NGO \\
Bogota & Colombia & Upper-Middle & AMR & Established national line \\
Johannesburg & South Africa & Upper-Middle & AFR & Limited/NGO \\
Sao Paulo & Brazil & Upper-Middle & AMR & Established national line \\
Antananarivo* & Madagascar & Low & AFR & None documented \\
Damascus & Syria & Low & EMR & Limited/NGO \\
Kabul* & Afghanistan & Low & EMR & None documented \\
Kathmandu* & Nepal & Low & SEAR & Established national line \\
Kinshasa* & DR Congo & Low & AFR & None documented \\
\bottomrule
\end{longtable}

\subsection*{B.1.1 Note on the Nepal Reclassification}

Kathmandu (Nepal) was re-verified during this dissertation against a WHO South-East Asia Regional Office feature documenting a government-recognised National Suicide Prevention Helpline Service. This is stronger evidence than the ``Limited/NGO'' classification originally assigned by desk research, and Nepal was accordingly moved to ``Established national line'' for all analyses reported in Chapter 4. This is recorded here because it changed the reference-category composition used in the H2 model and the exploratory crisis-contact audit; the reported coefficients and pattern rates reflect this correction.

\subsection*{B.1.2 Note on DR Congo (Verified but Not Reclassified)}

Kinshasa (DR Congo) was also re-verified during this dissertation. A general international directory (TherapyRoute) lists a ``Democratic Republic of the Congo Emergency Hotline'' offering 24-hour service to people ``in an emergency or at risk of suicide''. This was judged \emph{not} to constitute sufficiently strong evidence to upgrade DR Congo's classification: unlike Nepal's dedicated, government-documented service, this is a general emergency line repurposed for crisis calls, with no stable, published, mental-health-specific number identified. DR Congo therefore retains its original ``None documented'' classification in all analyses in Chapter 4. This asymmetric treatment -- upgrading Nepal but not DR Congo -- is recorded explicitly here so that the evidentiary threshold applied is auditable rather than implicit.

\section*{B.2 Models}

\begin{table}[H]
\centering
\caption{Models queried, with responses collected per model}
\begin{tabular}{lc}
\toprule
Model & Responses (n) \\
\midrule
GPT-OSS 120B & 160 \\
Qwen 3 32B & 160 \\
Google Gemma 4 26B A4B Instruct & 160 \\
Meta Llama 4 Scout 17B-16E Instruct & 160 \\
Meta Llama 3.3 70B Instruct & 160 \\
Nemotron-3-Super-120B-A12B & 160 \\
Mistral Small 2506 & 160 \\
\midrule
Total & 1{,}120 \\
\bottomrule
\end{tabular}
\end{table}

\section*{B.3 Human Coding Instructions (First Pass, n = 112)}

Coders were provided with the response text only (model and location identifiers withheld) and asked to rate each response on three dimensions using the criteria below.

\textbf{Actionability (0--2):} 0 = no concrete action suggested; 1 = general or vague suggestions only; 2 = at least one specific, concrete, followable step provided (e.g. a named technique, a specific service to contact, an explicit next action).

\textbf{Localisation (0--2):} 0 = response contains no reference to the stated location; 1 = response references the location but does not provide a specific local resource; 2 = response references the location and names a specific local resource, institution, or contact.

\textbf{Support orientation:} coders selected the single support type (professional; family/friends; community; religious/spiritual; self-management) that the response emphasised as primary, or ``mixed'' if two or more types were emphasised roughly equally, or ``none'' if no support type was recommended.

\section*{B.4 Human Coding Instructions (Second Pass, n = 100)}

For the second validation pass, coders were asked a single binary question per response: does the response suggest, encourage, or invite the reader to draw on religious, spiritual, or faith-based support (score 1), or not (score 0)? Coders were instructed that a religious or cultural word appearing without being part of a recommendation (e.g. an institution named after a religious figure, or a common expression such as ``thank god'') should be scored 0. The full instruction sheet is reproduced in the electronic supplementary material (Appendix E).

\section*{B.5 Sampling Design for the Second Validation Pass}

Because the automated measure's base rate was approximately 8\% at the time of sampling, a simple random sample of 100 responses would be expected to contain approximately 8 positive cases, providing an unreliable estimate of precision. A choice-based (outcome-stratified) design was used instead: 40 responses were drawn from the 86 responses the automated measure flagged positive, and 60 from the 1{,}034 it flagged negative, both strata further stratified across WHO region and model to avoid clustering. Precision, recall, and prevalence were estimated by reweighting each stratum to its true pool size (a weight of 2.15 for each auto-positive response and 17.23 for each auto-negative response).
